% Part: incompleteness
% Chapter: arithmetization-syntax
% Section: introduction

\documentclass[../../../include/open-logic-section]{subfiles}

\begin{document}

\olfileid{inc}{art}{int}
\olsection{پېژندګلو}

د محاسبه‌پوهې او منطق د نښلولو دپاره داسې لار پکار ده چې د منطق د
څيزونو (سمبولونو، ترمونو، !!{formula}s او !!{derivation}s)، پر هغو د
عملياتو، او د هغو د خاصيتونو او اړيکو په اړه داسې خبرې پرې وکړو چې
محاسبوي چلند ته برابرې وي۔ دا کار نېغ په نېغه کولے شو: پر سمبولونو،
د سمبولونو پر لړيو او له هغو څخه پر جوړو شويو نورو څيزونو محاسبه
کېدونکې تابعې او اړيکې څېړو۔ ځکه چې د منطقي نحو ټول څيزونه متناهي دي
او د سمبولونو له !!a{enumerable} سټ څخه جوړ شوي، نو د محاسبې په ځينو
مدلونو کښې دا شوني دي۔ خو د محاسبې نور مدلونه---لکه بازګشتي تابعې---
يوازې عددونو او د هغو اړيکو او تابعو ته محدود دي۔ برسېره پر دې، په
پاې کښې غواړو چې نحو د ځينو تيوريو دننه هم وڅېړو، په ځانګړي ډول په
هغو تيوريو کښې چې د حساب په ژبه کښې بيان شوې وي۔ په دې حالتونو کښې
د نحو \emph{حسابي کول} اړين دي؛ يعنې نحوي څيزونه، پر هغو عمليات او د
هغو اړيکې په ترتيب د عددونو، حسابي تابعو او حسابي اړيکو په توګه
تمثيل کړو۔ دا مفکوره چې لايبنېڅ ته رسېږي، نحوي څيزونو ته د عددونو
ټاکل دي۔

سمبولونو ته د هغو د «کوډونو» په توګه د عددونو ټاکل نسبتاً ساده دي۔
ځينې سمبولونه لږه ستونزه پيدا کوي، ځکه د بېلګې په توګه بې‌شمېره
!!{variable}s، او ان د هرې ځاييزې کچې~$n$ بې‌شمېره !!{function}s شته۔ خو
البته سمبولونو ته په منظم ډول داسې عددونه ټاکل کېدے شي چې، د بېلګې
په توګه، $\Obj v_2$ او $\Obj v_3$ بېل کوډونه ولري۔ د سمبولونو لړۍ
(لکه ترمونه او !!{formula}s) تر دې لويه ستونزه ده۔ خو که د عددونو له
لړيو سره سوچه حسابي چلند کولے شو (د بېلګې په توګه د لړيو د اولي
عددونو-د-قوتونو کوډول)، نو د جلا سمبولونو کوډول د سمبولونو د لړيو تر
کوډولو غځولے شو، او بيا ئې د !!{formula}s تر لړيو يا نورو ترتيبونو،
لکه !!{derivation}s، نور هم غځولے شو۔ دا غځېدلے کوډول «ګوډل شمېرنه»
بلل کېږي۔ هر ترم، !!{formula} او !!{derivation} ته يوه ګوډل شمېره
ټاکل کېږي۔

کله چې د سمبولونو لړۍ د هغو د کوډونو د لړۍ په توګه کوډوو، او د لړيو
د کوډولو داسې نظام غوره کوو چې په محاسبه کېدونکو تابعو پرې چلند کېدے
شي، نو پر ګوډل شمېرو هم په محاسبه کېدونکو تابعو چلند کولے شو۔ په عمل
کښې به ټولې اړوندې تابعې بنسټيزې بازګشتي وي۔ د بېلګې په توګه، د يوې
لړۍ اوږدوالے محاسبه کول او د لړۍ له کوډ څخه د هغې $i$-م غړے محاسبه
کول دواړه بنسټيز بازګشتي دي۔ که د لړۍ کوډوونکے عدد، د بېلګې په توګه،
د !!a{formula}~$!A$ ګوډل شمېره وي، سمدستي وينو چې د !!a{formula}
اوږدوالے او په !!a{formula} کښې د $i$-م سمبول (کوډ) هم د~$!A$ له
ګوډل شمېرې څخه محاسبه کېدے شي۔ د دې ثابتول لږ ستونزمن دي چې، د بېلګې
په توګه، د سم جوړ شوي ترم يا د سم !!{derivation} د ګوډل شمېرې په
توګه د يو عدد خاصيت بنسټيز بازګشتي دے۔ بيا هم دا شوني دي، ځکه د پام
وړ لړۍ (ترمونه، !!{formula}s او !!{derivation}s) په استقرايي ډول
تعريف شوې دي۔

د بېلګې په توګه د تعويض عمليه وګورئ۔ که $!A$ يو فارمول، $x$ يو
متغير او $t$ يو ترم وي، نو $\Subst{!A}{t}{x}$ په~$!A$ کښې د~$x$ د
هرې ازادې پېښې په~$t$ بدلولو پايله ده۔ اوس فرض کړئ چې $!A$، $x$ او
$t$ ته مو په ترتيب $k$، $l$ او $m$ ګوډل شمېرې ټاکلې دي۔ هماغه
طرحه $\Subst{!A}{t}{x}$ ته، مثلاً، ګوډل شمېره~$n$ ټاکي۔ دا تطبيق---
چې $k$، $l$ او $m$ پر~$n$ تطبيقوي---د تعويض د عمليې حسابي سيال دے۔
کله چې د تعويض عمليه $!A$، $x$ او $t$ پر
$\Subst{!A}{t}{x}$ تطبيقوي، د حسابي شوي تعويض تابع د ګوډل شمېرو
$k$، $l$ او $m$ پر ګوډل شمېرې~$n$ تطبيقوي۔ وبه وينو چې دا تابع
بنسټيزه بازګشتي ده۔

د نحو حسابي کول يوازې انتزاعي ارزښت نۀ لري، که څه هم په لومړي ځل دا
يوه نااشنا ژوره مفکوره وه چې د حساب د ژبې په شان ژبې، چې د ژبو «په
اړه د خبرو» کومه اله ورسره نۀ وي، بيا هم د عباراتو پېچلي خاصيتونه
صوري کولے شي۔ له دې وروسته يوازې يو وړوکے ګام پاتې کېږي چې وپوښتو په
دې ژبه کښې يوه تيوري، لکه د پېانو حساب، د خپلې ژبې په اړه څه
\emph{ثابتولے} شي (د بېلګې په توګه دا چې !!{sentence}s د ثبوت وړ يا
رښتيا دي که نۀ)۔ دا موږ د ګوډل (د ناثبوت‌وړتيا په اړه) او تارسکي (د
رښتيا د ناتعريفېدنې په اړه) نومياليو محدودوونکو قضيو ته رسوي۔ خو د
نحو د حسابي کولو چل د محاسبه‌پوهې د ځينو مهمو پايلو د ثبوت دپاره هم
مهم دے، د بېلګې په توګه د تيوريو د محاسبوي ځواک يا د محاسبې د بېلو
مدلونو ترمنځ د اړيکې په اړه۔ د نحو حسابي کول د نورو څيزونو او
خاصيتونو د حسابي کولو د يو مدل په توګه کار کوي۔ د بېلګې په توګه، په
ورته ډول د حالتونو او محاسبو (مثلاً د ټيورينګ ماشينونو د محاسبو)
حسابي کول هم شوني دي۔ په دې سره د محاسبې په يو مدل (لکه ټيورينګ
ماشينونو) کښې محاسبې په بل مدل (لکه بازګشتي تابعو) کښې شبيه کولے شو۔

\end{document}
